\subsection{What limits the resolution}
\label{sec:limites}

In the nominal draw of the test set, the RMSE in $z$ of the calibrated barycentre falls from \num{55.18} to \qty{37.02}{\cm} between efficiencies $\varepsilon$ of \num{0.005} and \num{1}, and from \num{42.45} to \qty{35.24}{\cm} between energies of \num{10} and \qty{100}{\MeV} (Figure~\ref{fig:baricentro_luz} shows means). With \qty{177.8}{\pe\per\MeV} on average, near the average yield of APEX, it stands at \qty{38.23}{\cm} and, below $z = {}$\qty{20}{\m}, with less light, at \qty{13.79}{\cm}, close to the \qty{13.7}{\cm} obtained by APEX at \qty{10}{\MeV} (Table~\ref{tab:referencias}). Detecting every photon would reduce the error of the barycentre by \qty{5.43}{\cm} (\qty{4.96}{\cm} averaged over the draws), and by more still without the rounding of Equation~\eqref{eq:luz}, which at \qty{10}{\MeV} sets the weakest channels to zero and raises this error; for this reason, at the same detected light, the energy sweep lies below the efficiency sweep. Over the whole test set, the error of the barycentre is mainly systematic: with $\varepsilon = {}$\num{0.02}, GATv2-cascade (\qty{20.46}{\cm}) already errs less in $z$ than the barycentre with every photon. This systematic part is due to the vertices from \qty{20}{\m} upwards, where the barycentre stays at \qty{63.16}{\cm} even with every photon (Section~\ref{sec:onde_erro}); below \qty{20}{\m}, its error is mainly statistical: from \qty{23.29}{\cm}, near GATv2-cascade (\qty{22.54}{\cm}), it falls to \qty{8.33}{\cm} with every photon.

The error of GATv2-cascade also falls with more light (Section~\ref{sec:onde_erro}), but the light decreases away from the long walls, and the models were trained only at \qty{10}{\MeV} and $\varepsilon = {}$\num{0.02}: for them, no sweep separates the effect of photon statistics from that of position. The influence of the grid coverage was not isolated. With the full training set, every test vertex has a training vertex at most \qty{50}{\cm} away; with a quarter of the events, as in the first comparison (Section~\ref{sec:dificuldades}), \qty{27.6}{\percent} of the test vertices lack such a neighbour, and the mean distance rises from \num{49.31} to \qty{56.18}{\cm}. The models also read \num{33} hand-crafted event-level scalars, and, in the final configuration, no experiment separated the contribution of message passing from that of the scalars: the result is that of the graph encoders with these scalars, not of message passing or of the graph representation alone.
